\section{Results}
\label{sec:results}

Unless stated otherwise, results are for the main operating condition and the fault-frequency features, calibrated
results use temperature scaling, and
values are bootstrap means with 95\,\% intervals in brackets (Section~\ref{sec:bootstrap}); values given as medians or
means over rotations are labelled as such.

% Generated by scripts/build_paper_assets.py -- do not edit by hand.
\begin{table*}[t]
\centering
\small
\begin{tabular}{lccc}
\toprule
Metric & \textbf{Random Forest} & \textbf{SVM} & \textbf{XGBoost} \\
\midrule
In-domain accuracy & 0.91 & 0.88 & 0.91 \\
In-domain ECE & 0.01 & 0.09 & 0.06 \\
Real accuracy & 0.69 [0.52, 0.86] & 0.69 [0.50, 0.86] & 0.70 [0.52, 0.85] \\
Real ECE & 0.19 [0.05, 0.35] & 0.17 [0.04, 0.35] & 0.20 [0.06, 0.35] \\
AURC & 0.15 [0.04, 0.34] & 0.15 [0.03, 0.34] & 0.15 [0.04, 0.33] \\
Error at 5\% target & 0.27 [0.10, 0.46] & 0.24 [0.07, 0.45] & 0.26 [0.09, 0.44] \\
Automated at 5\% target & 0.87 [0.66, 1.00] & 0.82 [0.47, 1.00] & 0.86 [0.69, 1.00] \\
Best possible automation at 5\% target & 0.32 [0.00, 0.71] & 0.31 [0.03, 0.75] & 0.30 [0.00, 0.72] \\
Conformal coverage (nominal 0.90) & 0.65 [0.43, 0.84] & 0.68 [0.45, 0.87] & 0.67 [0.45, 0.85] \\
\bottomrule
\end{tabular}
\caption{Main results (fault-frequency features, temperature scaling). In-domain: leave one source bearing out, raw probabilities. Real damage: mean [95\,\% interval] from 2,000 bootstrap draws over calibration rotations and test bearings.}
\label{tab:main}
\end{table*}


\subsection{The artificial-to-real gap}

On unseen artificial bearings, all three classifiers generalise well: leave-one-bearing-out accuracy is 0.91
(Random Forest), 0.88 (SVM) and 0.91 (XGBoost), and their probabilities are close to calibrated (ECE 0.01, 0.09 and
0.06; Table~\ref{tab:main}). On real damage, accuracy falls to 0.69--0.70 for every model (e.g.\ Random Forest
0.69 [0.52, 0.86]), and ECE rises to 0.17--0.20 after temperature scaling. The wide intervals reflect how strongly
individual real bearings differ (Section~\ref{sec:silent}).

Recall on real damage differs by class: 0.82--0.85 for healthy bearings, 0.73 for inner-race and 0.56--0.57 for
outer-race damage (means over rotations). Outer-race damage, the easiest class on artificial bearings (0.97--0.99
in-domain), is the hardest on real ones.

\subsection{Confidence thresholds chosen on artificial damage do not transfer}

Fig.~\ref{fig:riskcoverage} shows the error among automatically decided cases as more cases are automated. On
unseen artificial bearings, the best possible threshold keeps the error at or below 5\,\% while automating 81--91\,\% of cases.
On real damage, the best possible threshold keeps the error at or below 5\,\% for only 25--30\,\% of cases (median over
rotations, uncalibrated probabilities; the bootstrap means after temperature scaling in Table~\ref{tab:main} are
0.30--0.32).

\begin{figure*}[t]
  \centering
  \includegraphics[width=\textwidth]{fig1_risk_coverage}
  \caption{Error among automatically decided cases as a function of the share of cases automated (most confident
    first), for unseen artificial bearings (leave one bearing out; solid) and real damage (dashed; median and
    interquartile range over 105 rotations; uncalibrated probabilities). Dotted line: 5\,\% error target. With the best possible threshold, the
    models automate 81--91\,\% of unseen artificial-bearing cases at or below 5\,\% error, but only 25--30\,\% of
    real-damage cases.}
  \label{fig:riskcoverage}
\end{figure*}

In deployment the threshold must be chosen without real-damage labels. A threshold chosen on the calibration set for
a 5\,\% error target automates 82--87\,\% of real-damage cases, with an error of 0.24--0.27 among them
(Table~\ref{tab:main}). The lower ends of the 95\,\% intervals (0.07--0.10) are still above the target, and the
target is exceeded in all 105 rotations for every model. A stricter 1\,\% target does not solve the problem: about
half of the cases are still automated (0.48--0.53), with an error of 0.13--0.15 (means over rotations).

The reliability diagrams in Fig.~\ref{fig:reliability} show why: on real damage, observed accuracy falls below the
predicted confidence across almost the whole range---the models become systematically overconfident. Post-hoc
calibration on artificial bearings does not correct this reliably. Temperature scaling lowers ECE on real damage for
the SVM (0.19 $\to$ 0.17) and XGBoost (0.25 $\to$ 0.20) but raises it slightly for Random Forest (0.18 $\to$ 0.19);
isotonic regression gives 0.23--0.24 for all three, worse than temperature scaling for every model and worse than
no calibration for Random Forest and the SVM.

\begin{figure*}[t]
  \centering
  \includegraphics[width=\textwidth]{fig2_reliability}
  \caption{Reliability diagrams (raw probabilities) for unseen artificial bearings (leave one bearing out) and for
    real damage (pooled over 105 rotations). Bins with fewer than 20 windows are omitted. Dotted line: perfect
    calibration. Models that are close to calibrated on unseen artificial bearings are overconfident on real damage
    (ECE 0.18--0.25). The SVM's in-domain dip at 0.75 confidence comes from one held-out artificial bearing (KI01),
    which contributes 80 of that bin's 131 windows, all misclassified.}
  \label{fig:reliability}
\end{figure*}

\subsection{Conformal prediction under-covers}

Split conformal prediction with a nominal coverage of 0.90 reaches 0.65--0.68 on real damage (Table~\ref{tab:main}),
and even the upper ends of the intervals (0.84--0.87) stay below nominal. Mean prediction-set sizes are below
one (0.90--0.96 classes, means over rotations), so some sets are empty: 83--89\,\% of sets contain a single class, and 25--27\,\% of those single-class sets are wrong.
Class-conditional (Mondrian) conformal prediction is no better (coverage 0.63--0.64, means over rotations).

\subsection{Silent bearings}
\label{sec:silent}

Fig.~\ref{fig:perbearing} shows that the failures concentrate on a few bearings. Two outer-race bearings, KA15 and
KA22, are almost never recognised (accuracy 0.02--0.07), and two inner-race bearings, KI16 and KI17, are recognised in
only 29--44\,\% of windows. The models remain confident on all four: mean confidence is 0.69--0.94, and 0.83--0.93
for the two outer-race bearings. The remaining ten test bearings are recognised in at least 68\,\% of windows.
Confidently misclassified bearings are rare among unseen source bearings---none for Random Forest and XGBoost,
one for the SVM (KI01, accuracy 0.00 at mean confidence 0.76)---so calibration fitted on artificial damage has little
chance to anticipate them.

\begin{figure*}[t]
  \centering
  \includegraphics[width=\textwidth]{fig3_per_bearing}
  \caption{Mean confidence and accuracy of each test bearing (11 with real damage, 3 healthy), averaged over 105 rotations (raw
    probabilities). Dotted line: confidence equals accuracy. Four bearings (KA15, KA22: outer race; KI16, KI17: inner
    race) are mostly misclassified while the models remain 69--94\,\% confident on average.}
  \label{fig:perbearing}
\end{figure*}

\subsection{Feature ablation}
\label{sec:ablation}

Adding the five amplitude-based time-domain features (Table~\ref{tab:ablation}) lowers in-domain accuracy from
0.88--0.91 to 0.49--0.77, while accuracy on real damage barely changes (0.65--0.69 vs 0.69). The amplitude features
describe individual bearings rather than damage: unseen healthy bearings in particular are rarely recognised with
them (in-domain healthy accuracy 0.00--0.28), in line with the kurtosis differences between healthy bearings noted in
Section~\ref{sec:splits}. With all features, the in-domain reference is itself poor: for Random Forest and the SVM,
accuracy on real damage is even higher than on unseen artificial bearings (0.69 vs 0.62 and 0.65 vs 0.49), which
hides the artificial-to-real gap. The trustworthiness failure remains: the 5\,\% target is exceeded in
59--92\,\% of the rotations that automate any case, with errors of 0.12--0.17 (means over rotations, as in
Table~\ref{tab:ablation}). Conformal coverage is higher with all features (0.80--0.89), because the prediction sets are
larger (1.3--1.9 classes on average, against 0.9--1.0 with fault-frequency features only).

% Generated by scripts/build_paper_assets.py -- do not edit by hand.
\begin{table}[t]
\centering
\small
\caption{Feature ablation: all nine features vs fault-frequency features only. Means over 105 calibration rotations; temperature scaling; marginal split conformal prediction, nominal coverage 0.90.}
\label{tab:ablation}
\resizebox{\textwidth}{!}{%
\begin{tabular}{llccccccc}
\toprule
\makecell[l]{Features} & Model & \makecell{In-domain\\accuracy} & \makecell{In-domain\\ECE} & \makecell{Real\\accuracy} & \makecell{Real\\ECE} & \makecell{Error at\\5\% target} & \makecell{Runs exceeding\\5\% target} & \makecell{Conformal\\coverage} \\
\midrule
All 9 features & Random Forest & 0.62 & 0.12 & 0.69 & 0.16 & 0.15 & 92\% & 0.85 \\
All 9 features & SVM & 0.49 & 0.37 & 0.65 & 0.16 & 0.12 & 59\% & 0.89 \\
All 9 features & XGBoost & 0.77 & 0.14 & 0.68 & 0.16 & 0.17 & 91\% & 0.8 \\
Fault-frequency only & Random Forest & 0.91 & 0.01 & 0.69 & 0.19 & 0.27 & 100\% & 0.65 \\
Fault-frequency only & SVM & 0.88 & 0.09 & 0.69 & 0.17 & 0.24 & 100\% & 0.68 \\
Fault-frequency only & XGBoost & 0.91 & 0.06 & 0.69 & 0.2 & 0.26 & 100\% & 0.66 \\
\bottomrule
\end{tabular}}
\end{table}


\subsection{1D-CNN baseline}
\label{sec:cnnresults}

The 1D-CNN trained on raw vibration does not generalise to unseen bearings even within the source domain
(Table~\ref{tab:main}). Its leave-one-bearing-out accuracy is 0.46, and its errors are confident (ECE 0.45): every
held-out healthy bearing is classified as damaged at a mean confidence of 0.94--1.00, and four of the five artificial
inner-race bearings are misclassified. The network does learn damage-related patterns---it recognises six of the
seven artificial outer-race bearings, and the real-damage bearings KI14 and KI18 in 93\,\% and 100\,\% of
windows---but with 12 to 14 training bearings it also learns the signatures of individual bearings, in line with the
observation that the number of training bearings limits generalisation \citep{vieira2026}. The random-split sanity
check confirms this. When the test windows come from bearings that are also in the training set, the CNN classifies
98\,\% of them correctly with an ECE of 0.02, as accurately as the feature-based classifiers (0.97--0.98). Holding
out whole bearings costs the feature-based classifiers 6--10 percentage points of accuracy, and the CNN 51 points.

On real damage, the CNN's accuracy is 0.43 [0.16, 0.71]. Because its calibration bearings expose the same problem,
post-hoc calibration makes it cautious rather than reliable: temperature scaling lowers ECE from 0.44 to 0.19, and the
threshold for a 5\,\% error target automates only 13\,\% of cases, yet these still have an error of 0.38. In 55 of the
105 rotations the threshold automates no case at all; in the other 50 the target is exceeded in 43. Even the best
possible threshold automates only 10\,\% of cases.
Marginal conformal prediction reaches a coverage of 0.94, but only by returning 2.5 of the three classes on average;
class-conditional coverage is 0.54. Calibrating on one or two labelled real bearings per class did not help reliably
(error at the 5\,\% target 0.22--0.48 across calibration sets; target met in at most 49\,\% of draws). The CNN's
errors are also of a different kind: only 3\,\% of them are missed faults, the rest being false alarms on healthy
bearings and confusions between fault types. Of the four silent bearings, KA22, KI16 and KI17 remain misclassified
(accuracy 0.02, 0.27 and 0.11), while KA15 is recognised in 48\,\% of windows.

With AdaBN, the network's accuracy on real damage rises from 0.43 to 0.51 [0.24, 0.77], but its confidence does not
become more trustworthy (Table~\ref{tab:main}). The calibration error after temperature scaling is practically
unchanged (0.20 against 0.19), and the error at the 5\,\% target is 0.42. Because the calibration bearings, and with
them the thresholds, are the same as without adaptation, the threshold again automates no case in 55 rotations; in
all of the other 50 it now exceeds the target. Marginal conformal coverage is 0.89, again with large sets (2.3 classes
on average). The gain comes from outer-race and healthy bearings: KA15 and KA30, for example, are recognised in
90\,\% and 96\,\% of windows instead of 48\,\% and 59\,\%. Real inner-race damage, by contrast, is now almost always
called outer-race damage, and more confidently than before: KI04, KI16 and KI17 are classified correctly in 0--6\,\%
of windows at a mean confidence of 0.93--0.97.

\subsection{Calibrating on labelled real bearings}

Table~\ref{tab:realcal} and Fig.~\ref{fig:realcal} compare calibration sets on identical test bearings. With source
(artificial) calibration, the error at the 5\,\% target is 0.26--0.28 and the target is met in at most 1\,\% of draws.
Calibrating on one real bearing per class lowers the error to 0.16--0.19 and meets the target in 29--42\,\% of draws;
two real bearings per class lower it to 0.10--0.11 and meet it in 49--52\,\% of draws. Combining the artificial and
the real calibration bearings lies in between: the error at the 5\,\% target is 0.18--0.19 with one and 0.13--0.14
with two real bearings per class, with 0.58--0.62 and 0.44--0.46 of cases automated, and the target is met in at most
33\,\% of draws.

The improvement comes mainly from caution: the share of cases automated falls from 0.81--0.88 to 0.51--0.52 (one
bearing per class) and 0.35--0.38 (two). Calibration error improves moderately (ECE 0.19--0.22 $\to$ 0.13--0.17 with
two bearings per class), and conformal coverage rises from 0.64--0.67 to 0.84--0.85, still below nominal. Variation
across draws is large (Table~\ref{tab:realcal}), as it depends on which real bearings are available for calibration.

\begin{figure*}[t]
  \centering
  \includegraphics[width=\textwidth]{fig4_real_calibration}
  \caption{Calibrating on a few labelled real-damage bearings (temperature scaling; confidence threshold chosen on
    the calibration set for 5\,\% error). (a) Error among automated real-damage cases; (b) share of real-damage cases
    automated. Points: mean over draws; bars: 2.5th--97.5th percentile over draws (50 draws per real-bearing setting;
    artificial values pooled over all 100).}
  \label{fig:realcal}
\end{figure*}

% Generated by scripts/build_paper_assets.py -- do not edit by hand.
\begin{table*}[t]
\centering
\small
\resizebox{\textwidth}{!}{%
\begin{tabular}{llccccc}
\toprule
Model & Calibrated on & \makecell{Error at\\5\% target} & \makecell{Draws meeting\\5\% target} & Automated & ECE & \makecell{Conformal\\coverage} \\
\midrule
Random Forest & Artificial bearings & 0.28 [0.13, 0.42] & 0\% & 0.87 [0.65, 1.00] & 0.21 [0.10, 0.33] & 0.64 [0.49, 0.80] \\
Random Forest & 1 real bearing/class & 0.19 [0.00, 0.36] & 29\% & 0.51 [0.00, 0.99] & 0.20 [0.05, 0.32] & 0.78 [0.52, 0.98] \\
Random Forest & 2 real bearings/class & 0.11 [0.01, 0.38] & 51\% & 0.38 [0.16, 0.77] & 0.17 [0.02, 0.33] & 0.85 [0.57, 0.99] \\
SVM & Artificial bearings & 0.26 [0.12, 0.39] & 1\% & 0.81 [0.45, 1.00] & 0.19 [0.04, 0.33] & 0.67 [0.51, 0.84] \\
SVM & 1 real bearing/class & 0.18 [0.00, 0.45] & 35\% & 0.52 [0.00, 0.99] & 0.18 [0.04, 0.34] & 0.77 [0.45, 0.98] \\
SVM & 2 real bearings/class & 0.10 [0.00, 0.36] & 52\% & 0.37 [0.18, 0.76] & 0.13 [0.04, 0.29] & 0.85 [0.60, 1.00] \\
XGBoost & Artificial bearings & 0.28 [0.15, 0.41] & 0\% & 0.88 [0.70, 1.00] & 0.22 [0.11, 0.32] & 0.65 [0.51, 0.80] \\
XGBoost & 1 real bearing/class & 0.16 [0.00, 0.37] & 42\% & 0.52 [0.08, 0.99] & 0.19 [0.07, 0.33] & 0.79 [0.54, 0.97] \\
XGBoost & 2 real bearings/class & 0.11 [0.00, 0.37] & 49\% & 0.35 [0.03, 0.76] & 0.16 [0.06, 0.31] & 0.84 [0.58, 0.99] \\
\bottomrule
\end{tabular}}
\caption{Calibrating on artificial bearings vs one or two labelled real bearings per class (fault-frequency features, temperature scaling). Mean [2.5th, 97.5th percentile] over 50 draws per setting; artificial-bearing values pooled over all 100 draws. Within each draw, all settings are scored on the same test bearings. Draws meeting the target count only draws that automated at least one case.}
\label{tab:realcal}
\end{table*}


\subsection{Other operating conditions}
\label{sec:conditions}

Table~\ref{tab:conditions} repeats the evaluation at the three other operating conditions, with the same bearing
splits. At lower load torque the results are almost identical to the main condition: unseen artificial bearings are
classified with accuracy 0.81--0.94, real damage with 0.71, the 5\,\% target is exceeded in every rotation with
an error of 0.25--0.27, and conformal coverage is 0.64--0.68.

At lower radial force and lower speed, the fault-frequency features separate the classes less well even on artificial
damage (in-domain accuracy 0.70--0.76 and 0.57--0.69). The models are then less confident, and the failure is milder
but still present: the error at the 5\,\% target is 0.15--0.21 and 0.10--0.12, the target is exceeded in 88--99\,\% and
68--83\,\% of the rotations that automate any case, and conformal coverage reaches 0.80--0.86, with intervals that include the nominal 0.90.

KA15 and KA22 are misclassified at every condition (mean accuracy over the three models at most 0.41). KI17 is
misclassified at the main condition, at low torque and at 900\,rpm, and KI16 at the main condition and at 900\,rpm;
at low torque KI16 is recognised in 56\,\% of windows. The nature of the errors,
however, depends on the condition. Where the features are strong, most errors are missed faults (62--77\,\% of
errors at the main and low-torque conditions). At lower radial force and lower speed, false alarms on healthy bearings
also occur---healthy bearing K006, for example, is classified as damaged in 92\,\% of its windows at 400\,N---and missed
faults account for 20--27\,\% and 45--50\,\% of errors.

% Generated by scripts/build_paper_assets.py -- do not edit by hand.
\begin{table*}[t]
\centering
\small
\resizebox{\textwidth}{!}{%
\begin{tabular}{llcccccc}
\toprule
Condition & Model & \makecell{In-domain\\accuracy} & \makecell{Real\\accuracy} & \makecell{Error at\\5\% target} & \makecell{Runs exceeding\\5\% target} & \makecell{Conformal\\coverage} & \makecell{Missed faults\\(share of errors)} \\
\midrule
1500 rpm, 0.7 Nm, 1000 N (main) & Random Forest & 0.91 & 0.69 [0.52, 0.86] & 0.27 [0.10, 0.46] & 100\% & 0.65 [0.43, 0.84] & 0.67 \\
1500 rpm, 0.7 Nm, 1000 N (main) & SVM & 0.88 & 0.69 [0.50, 0.86] & 0.24 [0.07, 0.45] & 100\% & 0.68 [0.45, 0.87] & 0.62 \\
1500 rpm, 0.7 Nm, 1000 N (main) & XGBoost & 0.91 & 0.70 [0.52, 0.85] & 0.26 [0.09, 0.44] & 100\% & 0.67 [0.45, 0.85] & 0.68 \\
1500 rpm, 0.1 Nm, 1000 N & Random Forest & 0.94 & 0.71 [0.53, 0.87] & 0.27 [0.11, 0.46] & 100\% & 0.64 [0.41, 0.84] & 0.77 \\
1500 rpm, 0.1 Nm, 1000 N & SVM & 0.81 & 0.71 [0.52, 0.87] & 0.25 [0.08, 0.44] & 100\% & 0.68 [0.42, 0.89] & 0.69 \\
1500 rpm, 0.1 Nm, 1000 N & XGBoost & 0.94 & 0.71 [0.54, 0.87] & 0.26 [0.09, 0.44] & 100\% & 0.67 [0.43, 0.85] & 0.73 \\
1500 rpm, 0.7 Nm, 400 N & Random Forest & 0.76 & 0.66 [0.52, 0.82] & 0.17 [0.02, 0.39] & 98\% & 0.80 [0.56, 0.96] & 0.26 \\
1500 rpm, 0.7 Nm, 400 N & SVM & 0.70 & 0.65 [0.45, 0.81] & 0.21 [0.02, 0.46] & 99\% & 0.80 [0.56, 0.95] & 0.20 \\
1500 rpm, 0.7 Nm, 400 N & XGBoost & 0.74 & 0.67 [0.52, 0.82] & 0.15 [0.00, 0.36] & 88\% & 0.82 [0.59, 0.96] & 0.27 \\
900 rpm, 0.7 Nm, 1000 N & Random Forest & 0.69 & 0.61 [0.44, 0.78] & 0.11 [0.01, 0.33] & 79\% & 0.85 [0.65, 0.99] & 0.50 \\
900 rpm, 0.7 Nm, 1000 N & SVM & 0.57 & 0.61 [0.44, 0.78] & 0.10 [0.01, 0.30] & 68\% & 0.86 [0.64, 0.99] & 0.45 \\
900 rpm, 0.7 Nm, 1000 N & XGBoost & 0.68 & 0.61 [0.45, 0.77] & 0.12 [0.01, 0.33] & 83\% & 0.85 [0.66, 0.97] & 0.48 \\
\bottomrule
\end{tabular}}
\caption{All four Paderborn operating conditions (shaft speed, load torque, radial force), with the same bearing splits. Fault-frequency features, temperature scaling. Real damage: mean [95\,\% interval] from 2,000 bootstrap draws. Missed faults: share of real-damage errors that call a damaged bearing healthy.}
\label{tab:conditions}
\end{table*}



\subsection{Adaptive demodulation band}
\label{sec:adaptive}

Selecting the demodulation band with the fast kurtogram (Section~\ref{sec:features}) did not recover the silent
bearings; it lowered accuracy considerably. The kurtogram chose the band from 10.7 to 21.3\,kHz in 70\,\% of windows,
and this was the most frequent choice for 22 of the 29 bearings, all six healthy bearings included, so it reflects
impulsive content of the test rig rather than bearing damage. In the selected bands the fault peaks of most damaged
bearings weaken or disappear---the outer-race peak of artificial bearing KA07, for example, falls from 16.2 to 1.3
times the noise floor (medians over windows)---and a clear peak remains only for KA01 and for the two bearings whose
band lies near 10\,kHz (KI01 and KI18). The four silent bearings gain no peak (1.8--2.2 times the noise floor at
their fault frequency). Accuracy drops to 0.32--0.36 on unseen artificial bearings and to 0.41--0.43 on real damage
(main condition, temperature scaling). An earlier, simpler selection by spectral kurtosis with a fixed 4\,kHz band
width gave the same picture. The fixed 2--10\,kHz band is therefore the better choice for these signals. A band chosen
by another criterion might still reveal the silent bearings, but neither selection gives evidence for it.
